\section{Significance of the Problem}
\label{sec:significance}

\subsection{For users: a vocabulary they are never given}

To decide whether to install a cycle tracker, a person has to weigh a benefit
they can feel against a risk they have to imagine. Privacy communication is the
main material they have for the second half of that judgement, and it is
material about mechanisms: categories of data, categories of recipient,
retention periods. If it never names a consequence, the user must supply the
consequence themselves.

The evidence suggests they cannot. Fitness tracker users are confident about
what is collected and unsure how it is used, and they rate harmful scenarios as
plausible but unlikely~\parencite{gabriele2020understanding}. Period tracker
users report high concern coupled with a sense of being uninformed and powerless
about what to do~\parencite{cao2024deleted}. That pattern --- concern without
actionable understanding --- is what one would predict from communication that
is complete about data flows and silent about stakes. Establishing whether the
silence is real, and where exactly it falls, turns a plausible explanation into
a measured one.

The risk is sharpest where inference is sharpest. A user who reads that an app
does not share raw data may reasonably conclude they are protected. If the
provider derives and acts on an inference instead, the conclusion is wrong, and
nothing in the text told them so~\parencite{martin2025cookies}.

\subsection{For providers and designers: a map of where communication fails}

Providers are not, in the main, choosing between honesty and deception; the
measurement literature already covers deception. They are choosing what to put
in a finite amount of text and screen space, and the conventions they inherit ---
from legal templates, from store label schemas --- have no slot for harm. A
systematic map of which harms are named, on which surface, and in what terms
gives design and policy teams something they currently lack: an account of what
their communication omits, expressed in categories a non-lawyer can act on.

There is already evidence that the design space exists. Hudig and Singh
co-designed ten concrete transparency and control mechanisms with fertility
tracker users~\parencite{hudig2025intimate}, and Malki et al.\ derive concrete
recommendations from the inconsistencies they
found~\parencite{malki2024exploring}. Both stop short of asking what the
communication should \emph{say about harm}; this thesis supplies that input.

\subsection{For regulators: evidence about a gap the law leaves open}

Section~\ref{sec:law} set out the structural point: the GDPR's
transparency articles enumerate practices, and require a statement of
consequences only in the narrow automated-decision-making case of Article
13(2)(f); breach communication under Article 34 is
reactive~\parencite{gdpr2016}. Platform disclosure schemas inherit the same
frame, and are in any case unreliable~\parencite{arkalakis2024abandon,li2022ios}.

Whether harm-facing disclosure would add anything is currently an untested
assumption in that debate. If providers turn out to name harms voluntarily and
in useful terms, the case for mandating it weakens. If, as expected, the four
elements identified in Section~\ref{sec:problem} are near-absent across an
entire market segment handling Article~9 special-category data, that is a
concrete, citable finding about the limits of a transparency regime built on
data-practice disclosure.

The structured platform disclosures do not fill the gap either. Both are
developer self-declarations against a fixed schema of data categories, which is
why their measured error rates are as high as they
are~\parencite{arkalakis2024abandon,li2022ios} --- and, more to the point here,
why no schema of that shape can carry a statement about harm in the first place.
What providers say in the free text, where nothing constrains them, is therefore
the half of the picture that no existing measurement covers.

\subsection{For research: extending harm analysis to provider discourse}

The harm literature has been developed in two places: in legal scholarship,
where the object is what a court will
recognise~\parencite{citron2021harms,calo2011boundaries,solove2006taxonomy}, and
in systems analysis, where the object is what a system's functionality makes
possible~\parencite{knabenhans2026harms}. Neither has been turned on the text
that providers actually put in front of users.

Doing so makes three contributions to method as well as to findings. It tests
whether a typology built for adjudication can be operationalised as a coding
scheme for corporate prose --- a question with an answer either way. It extends
usable-privacy document analysis, whose annotation schemes are schemes of data
practices~\parencite{wilson2016creation,harkous2018polisis}, with a harm
dimension. And it joins document analysis to interface inspection in one design,
so that what a provider writes can be compared with what the same provider
shows~\parencite{wharton1994cognitive}, which is where RQ4 lives.
